% [4 + 4 pts] ~4 páginas. Dos estrategias con uso INTENSIVO de derivados: diseño, valorización y escenarios.
% Tablas y figuras SOLO desde generado/ (make calc) o con macros. Cada una con \fuente{...}.
\section{Estrategias de cobertura con instrumentos derivados}

\subsection{Estrategia 1: \emph{cross currency swap} y NDF para el riesgo cambiario}

\textbf{Diseño.} Se propone un CCS amortizable con un banco de primer nivel que replique el cronograma remanente del
préstamo del BN. PETROPERÚ \textbf{recibe} cada mes, en soles, la cuota del préstamo y
\textbf{paga} una cuota en dólares a tasa fija sobre un nocional de US\$~\CCSNocionalUSD\ millones. Así la deuda en
soles se convierte sintéticamente en deuda en la moneda funcional (Figura~\ref{fig:ccs}). El resto de la posición en
soles (S/~\RestoPEN\ millones, sobre todo cuentas por pagar) se cubre en un \NDFPct\ con NDF de compra de soles a
\NDFDias\ días, renovables, que se liquidan por diferencias en dólares sin entregar el principal. Esa posición usa
saldos de diciembre de 2025 (salvo el BN) y debe actualizarse antes de pactar; los swaps de corto plazo con Citibank
siguen vigentes (activo de US\$~\SwapCitiJunio\ millones a junio de 2026) sin que se revele su nocional, que debe
restarse del NDF si cubre soles. El EEFF no publica el cronograma del BN (\BNCuotasTot\ cuotas en 48 meses), pero a
junio de 2026 el saldo seguía en S/~\BNSaldoJunio\ millones y era íntegramente no corriente
\parencite{petroperu2026eeffjun}: no vence principal antes de julio de 2027. Se supone que se pagan solo
intereses hasta junio de 2027 (S/~\BNCuota\ millones al mes) y luego \BNCuotasAmort\ cuotas francesas de
S/~\BNCuotaAmort\ millones, con la TEA llevada a tasa mensual, $(1+\text{TEA})^{1/12}-1$; con otros cronogramas
plausibles la tasa del CCS varía entre \CCSTasaMin\ y \CCSTasaMax.

\begin{figure}[t]\centering\small
\begin{tikzpicture}[caja/.style={draw,rounded corners,align=center,minimum width=2.9cm,minimum height=0.9cm,font=\small\bfseries},
                    >={Stealth[length=2mm]},node distance=3.4cm]
  \node[caja,fill=orange!12] (bn) {Banco de la Nación\\(acreedor en S/)};
  \node[caja,fill=blue!8,right=of bn] (pp) {PETROPERÚ};
  \node[caja,fill=green!10,right=of pp] (bk) {Banco contraparte\\(CCS)};
  \draw[->] (pp) -- node[above,font=\footnotesize]{Cuotas en S/} node[below,font=\footnotesize]{\BNTasa, \BNCuotasRem\ meses} (bn);
  \draw[->] ([yshift=3mm]bk.west) -- node[above,font=\footnotesize]{Recibe las cuotas en S/} ([yshift=3mm]pp.east);
  \draw[->] ([yshift=-3mm]pp.east) -- node[below,font=\footnotesize]{Paga US\$ al \CCSTasaUSD} ([yshift=-3mm]bk.west);
\end{tikzpicture}
\caption{Estructura de flujos del CCS propuesto}\label{fig:ccs}
\fuente{\textcite{petroperu2026eeff}, Nota 14(ii), p.~74. Elaboración propia.}
\end{figure}

\textbf{Valorización.} La tasa fija en dólares del CCS iguala el valor presente de ambas patas al inicio. La pata en
soles se descuenta con la curva cupón cero soberana de la SBS (supuesto: no hay una curva swap en soles pública) y la
pata en dólares con la curva cero obtenida por \emph{bootstrap} de los swaps SOFR
\parencite{sbs2026curvas,bluegamma2026,hull2022}. Ambas patas se convierten al tipo de cambio de compra, porque
PETROPERÚ vende dólares; el registro contable usa el TC medio. Como las cuotas ya incluyen la amortización, no hay
intercambio final de nocional. El forward del NDF se obtiene por paridad cubierta de tasas, con el SOFR a tres meses
(tasa simple de \USDTresMesesSimple) convertido a tasa efectiva:
\begin{equation}
  \frac{1}{S_0}\sum_{k} C_k^{S/}\,DF^{S/}(t_k) \;=\; \sum_{k} C_k^{US\$}(r_{US\$})\,DF^{US\$}(t_k),
  \qquad
  F_{NDF} = S_0^{compra}\left(\frac{1+r_{S/}}{1+r_{US\$}}\right)^{d/360}.
\end{equation}

\begin{table}[t]\centering\small
\caption{E1: Parámetros y resultados al \FechaValorizacion}\label{tab:e1-param}
\newcommand{\grupo}[1]{\rowcolor{gray!12}\multicolumn{3}{l}{\textit{#1}}\\}
\begin{tabularx}{\linewidth}{>{\raggedright\arraybackslash}p{6.4cm} r >{\raggedright\arraybackslash}X}
\toprule
\textbf{Parámetro} & \textbf{Valor} & \textbf{Fuente / supuesto} \\
\midrule
\grupo{Insumos de mercado}
Tipo de cambio: compra / venta (S/ por US\$) & \SpotBid\ / \SpotAsk & BCRP, interbancario \\
Tasa cero en S/: 1 año / 3 años & \PENUnAnio\ / \PENTresAnios & SBS, curva soberana \\
Tasa swap en US\$: 1 año / 3 años & \USDUnAnio\ / \USDTresAnios & Swaps SOFR \\
\grupo{CCS amortizable: recibe S/, paga US\$}
Saldo cubierto del préstamo del BN (S/ MM) & \BNSaldoRem & Cronograma supuesto, \BNCuotasRem\ pagos \\
Nocional (US\$ MM) & \CCSNocionalUSD & Saldo al tipo de cambio de compra \\
VP de la pata en S/ (US\$ MM) & \CCSVPPataPEN & Descontada con la curva SBS \\
Tasa fija de equilibrio de la pata en US\$ & \CCSTasaUSD & Antes del cargo por riesgo de crédito \\
Tasa con CVA (\emph{all-in}) & \CCSTasaAllIn & CVA de \CCSCVApb~pb; LGD de \LGD \\
Pagos totales de la pata en US\$ (US\$ MM) & \CCSTotalUSD & Suma de los \BNCuotasRem\ pagos \\
\grupo{NDF de compra de S/ a \NDFDias\ días}
Nocional (S/ MM) & \NDFNocional & \NDFPct\ del resto de la posición en S/ \\
Tipo de cambio forward (S/ por US\$) & \NDFForward & \NDFPuntos\ pips sobre el spot de compra \\
\midrule
\textbf{Posición en S/ sin cubrir (S/ MM)} & \textbf{\ResiduoTotal} & Tras la estrategia: cobertura de \CoberturaFX \\
\bottomrule
\end{tabularx}
\fuente{\textcite{petroperu2026eeff}, Notas 3 y 14; \textcite{bcrpseries2026}; \textcite{sbs2026curvas}; \textcite{bluegamma2026}. Elaboración propia.}
\end{table}

La tasa de equilibrio de la pata en dólares (\CCSTasaUSD) no es comparable con el \BNTasa\ del préstamo, porque son
tasas en monedas distintas. Refleja dos efectos: en el plazo del CCS la curva swap en dólares (\USDDosAnios\ a dos
años) está por encima de la soberana en soles (\PENDosAnios), y el cupón del BN supera a esa curva soberana, de modo
que la pata en soles vale US\$~\CCSVPPataPEN\ millones frente a un nocional de US\$~\CCSNocionalUSD\ millones; esa
prima se traslada a la tasa en dólares. El costo de cubrirse lo mide el diferencial de tasas implícito en el forward.
El NDF se pacta a \NDFForward, \NDFPuntos~pips por debajo del spot de compra, porque la tasa en soles a tres meses
(\PENTresMeses, curva SBS a 91 días) es menor que la de dólares (\USDTresMeses). Como PETROPERÚ compra soles a plazo, recibe menos soles
por dólar: la cobertura le cuesta \NDFCostoAnual\ anual, unos US\$~\NDFCostoUSD\ millones por trimestre. Con la curva sintética en dólares de la SBS, implícita en los
forwards locales, el forward sería \NDFForwardSBS\ (\NDFDifSBSPips~pips de diferencia) y la tasa del CCS \CCSTasaSBS: el
basis local es pequeño, aunque la tasa ejecutable requiere cotizaciones bancarias. Aun así la
tasa del CCS queda muy por debajo del rendimiento que el mercado exigía a los bonos de la Compañía a diciembre de 2025 (\YTMBonos). Dada
la calificación B-/Caa1, el banco cobrará un ajuste por riesgo de crédito (CVA). Con la exposición esperada del banco
(máximo de US\$~\CCSEPEMax\ millones, volatilidad realizada del TC a un año de \VolTC), el spread de
\SpreadBonosPb~pb a diciembre de 2025 y una LGD de \LGD, el CVA es de US\$~\CCSCVAUSD\ millones, unos \CCSCVApb~pb
anuales: la tasa \emph{all-in} sería \CCSTasaAllIn\ \parencite{hull2022}. Con el spread de agosto de 2026
(\SpreadProxyPb~pb) el CVA bajaría a \CCSCVApbProxy~pb. Es un mínimo, porque no incluye el costo de fondeo del banco (FVA) ni el riesgo \emph{wrong-way}:
la exposición del banco crece si el sol se deprecia, lo que puede coincidir con un deterioro soberano o con un
crudo caro que presione la caja de la Compañía. El Cuadro~\ref{tab:e1-flujos} muestra
que la pata en soles calza con el servicio de deuda al BN, por lo que PETROPERÚ solo desembolsa dólares, su moneda
funcional.

% ARCHIVO GENERADO — NO EDITAR A MANO
\begin{table}[htbp]\centering\small
\caption{E1: Flujos del CCS por año calendario (millones)}\label{tab:e1-flujos}
\begin{tabular}{lrrrr}
\toprule
 & \textbf{Recibe S/} & \textbf{Interés S/} & \textbf{Paga US\$} & \textbf{Interés US\$} \\
\midrule
2026 & 51.0 & 51.0 & 15.7 & 15.7 \\
2027 & 1,411.6 & 190.2 & 413.9 & 58.4 \\
2028 & 2,619.4 & 75.2 & 763.4 & 23.1 \\
Total & 4,081.9 & 316.3 & 1,192.9 & 97.2 \\
\bottomrule
\end{tabular}
\fuente{PETROPERÚ, EEFF auditados al 31.12.2025 (Gaveglio, Aparicio y Asociados, 2026), Nota 14(ii); SBS, curva cupón cero soberana en soles (28-sep-2026); BlueGamma, swaps SOFR (28-sep-2026). Elaboración propia. La pata en S/ replica el servicio de deuda al BN; PETROPERÚ solo desembolsa la pata en US\$.}
\end{table}


\textbf{Posiciones y demostración de la cobertura.} A la fecha de valorización la posición de cambio contable (PCC)
es pasiva en S/~\PosicionHoy\ millones, igual que al cierre de 2025 porque el BN aún no amortiza. El CCS y el NDF crean una posición activa en derivados (PND), y la posición global (PCG) queda pasiva en solo
S/~\ResiduoTotal\ millones (Cuadro~\ref{tab:e1-posiciones}). En el NDF, lo que pierden las cuentas por pagar cuando el
sol se aprecia lo gana el derivado: el total es constante en todo escenario e igual a $N(1/S_0-1/F)$, el costo de la
cobertura. Por cada S/~1,000 de nocional, el NDF paga US\$~\NDFPorMilMenosDiez\ si el TC cae 10\,\% (Cuadro~\ref{tab:e1-ndf}).

\medskip\noindent
\begin{minipage}[t]{0.42\linewidth}\footnotesize\setlength{\tabcolsep}{2.5pt}% ARCHIVO GENERADO — NO EDITAR A MANO
\centering
\captionof{table}{E1: posiciones de cambio (S/ MM)}\label{tab:e1-posiciones}
\begin{tabular}{lrr}
\toprule
 & \textbf{S/ MM} & \textbf{TC −10\% (US\$)} \\
\midrule
PCC: pasivo neto en S/ & (4,704.0) & (152.0) \\
PND: CCS (recibe S/) & 3,765.6 & 121.7 \\
PND: NDF (compra S/) & 750.7 & 24.3 \\
PCG = PCC + PND & (187.7) & (6.1) \\
\bottomrule
\end{tabular}
\fuente{\textcite{petroperu2026eeff}, Notas 3 y 14. Elaboración propia. Negativo = posición pasiva en S/ (o pérdida). Efecto a TC medio; el Cuadro~\ref{tab:e1-ndf} usa el TC de compra pactado.}
\end{minipage}\hfill
\begin{minipage}[t]{0.56\linewidth}\footnotesize\setlength{\tabcolsep}{2.5pt}% ARCHIVO GENERADO — NO EDITAR A MANO
\centering
\captionof{table}{E1: NDF a 91 días, resultado al vencimiento (US\$ MM)}\label{tab:e1-ndf}
\begin{tabular}{lrrrr}
\toprule
 & \textbf{Inicio} & \textbf{TC -10\%} & \textbf{Spot} & \textbf{TC +10\%} \\
\midrule
Cuentas por pagar en S/ & 0.00 & (24.27) & 0.00 & 19.86 \\
NDF: N(1/S − 1/F) & 0.00 & 24.14 & (0.13) & (19.99) \\
Total & 0.00 & (0.13) & (0.13) & (0.13) \\
NDF por S/ 1,000 (US\$) & 0.00 & 32.16 & (0.17) & (26.63) \\
\bottomrule
\end{tabular}
\fuente{PETROPERÚ, EEFF auditados al 31.12.2025 (Gaveglio, Aparicio y Asociados, 2026), Nota 3; SBS, curva cupón cero soberana en soles (28-sep-2026); BlueGamma, swaps SOFR (28-sep-2026). Elaboración propia. TC compra (al que PETROPERÚ pacta); S de liquidación: 3.093 / 3.437 / 3.780. Inicio: valor cero, sin flujo. Total constante = N(1/S0 − 1/F).}
\end{minipage}
\medskip

% ARCHIVO GENERADO — NO EDITAR A MANO
\begin{table}[htbp]\centering\small
\caption{E1: Resultado por diferencia de cambio según TC final (US\$ millones)}\label{tab:e1-escenarios}
\begin{tabular}{lrrrrr}
\toprule
 & \textbf{3.094 (-10\%)} & \textbf{3.266 (-5\%)} & \textbf{3.438 (base)} & \textbf{3.610 (+5\%)} & \textbf{3.782 (+10\%)} \\
\midrule
Sin cobertura & (123.4) & (58.5) & 0.0 & 52.9 & 101.0 \\
Solo CCS & (30.3) & (14.4) & 0.0 & 13.0 & 24.8 \\
CCS + NDF & (6.1) & (2.9) & 0.0 & 2.6 & 5.0 \\
\bottomrule
\end{tabular}
\fuente{PETROPERÚ, EEFF auditados al 31.12.2025 (Gaveglio, Aparicio y Asociados, 2026), Notas 3 y 14; BCRP, series estadísticas (TC interbancario 28-sep-2026). Elaboración propia. TC medio (compra-venta). Negativo = pérdida. Posición en S/ estimada a la fecha de valorización.}
\end{table}


\begin{figure}[t]\centering
  \includegraphics[width=0.8\linewidth]{f_e1_escenarios}
  \caption{E1: Resultado cambiario con y sin cobertura según el TC al cierre}\label{fig:e1}
  \fuente{\textcite{petroperu2026eeff}, Notas 3 y 14; \textcite{bcrpseries2026}. Elaboración propia.}
\end{figure}

\textbf{Escenarios.} Con la posición a la fecha de valorización (S/~\PosicionHoy\ millones al TC medio de
\SpotUSDPEN), una caída de 10\,\% del tipo de cambio (una apreciación del sol de \ApreciacionSolDiez) costaría
US\$~\EscSinCobMenosDiez\ millones sin cobertura y solo US\$~\EscConCobMenosDiez\ millones con el CCS y los NDF: la
sensibilidad cae en \ReduccionFX\ (Cuadro~\ref{tab:e1-escenarios} y Figura~\ref{fig:e1}). La pérdida sin cobertura
difiere de los US\$~\SensApreciacion\ millones de la sección~2 porque aquella parte del TC de cierre de 2025. A cambio, la Compañía renuncia a la ganancia si el sol se deprecia; el objetivo es
estabilizar resultados y flujos, no especular \parencite{stulz1996}.

\textbf{Registro contable.} Ambos derivados se pactan a valor cero y luego se miden a valor razonable: como activo si
el valor es favorable y como pasivo si no. Ante una caída instantánea de 10\,\% del TC, el CCS sería un activo de
US\$~\CCSValMenosDiez\ millones y el NDF uno de US\$~\NDFValMenosDiez\ millones; si el TC sube 10\,\%, serían pasivos
de US\$~\CCSValMasDiez\ y \NDFValMasDiez\ millones. Difieren algo de la compensación nominal del
Cuadro~\ref{tab:e1-posiciones}, porque el cupón del BN supera la curva y el NDF aún no vence: esa brecha es la
inefectividad esperada. El CCS se designa como cobertura de flujos de efectivo del préstamo
del BN: la parte efectiva va a otro resultado integral y se reclasifica cuando la diferencia de cambio del préstamo
afecta resultados. La relación debe documentarse sobre el cronograma contractual; sus fuentes de inefectividad son las
diferencias de cronograma, el CVA y el basis, que puede excluirse de la designación (NIIF~9, 6.5.16). El NDF no se
designa: cubre partidas monetarias ya reconocidas, cuya diferencia de cambio va a resultados por la NIC~21, y medido a
valor razonable con cambios en resultados la compensa en el mismo estado, sin contabilidad de coberturas \parencite{ifrs9}.

\subsection{Estrategia 2: swap y collar de costo cero sobre WTI para el inventario de crudo}

\textbf{Diseño.} Lo que está expuesto es el desfase de precios: el crudo se paga al precio de su embarque y los
productos se venden por paridad de importación semanas después, así que el inventario pierde valor si el precio cae
antes de trasladarse a ventas. Las compras aún sin precio no suman exposición, porque las ventas que financiarán
fijarán su precio en el mismo mercado. Con compras de \ComprasMBDC\ MBDC, las \InvCrudoMbl\ MBL rotan en unos
\RotacionDias\ días. Por eso la cobertura se arma en \NCapas\ \textbf{capas mensuales}: un tercio del volumen en los
contratos de noviembre, diciembre y enero, cada capa sobre el inventario que se venderá ese mes. Se cubre el
\CoberturaCrudoPct\ (\InvCubiertoMMbl\ MMbl), tope de la política propuesta, del inventario al cierre de 2025, que al
WTI de US\$~\WTISpot/bl vale US\$~\InvValorHoy\ millones. No se cubre el inventario de productos
(US\$~\InvRefinadosUSD\ millones), cuyo precio ya sigue la paridad con menor desfase. Se usan instrumentos OTC, porque
los futuros NYMEX exigen márgenes diarios. El \RepartoSwap\ del volumen va a (a) un \textbf{swap de WTI} con
liquidaciones mensuales a US\$~\SwapPrecio/bl (promedio de los futuros de noviembre a enero), para cargamentos con
margen ya definido; el resto, a (b) un \textbf{collar de costo cero} por capa: un put comprado al \PutPctFwd\ del
futuro de cada mes (US\$~\CollarPutsCapas) financiado con un call vendido (US\$~\CollarCallsCapas). Las opciones
vencen en \PlazoDiasCapas\ días y se valorizan con el modelo de \textcite{black1976}:
\begin{equation}
  P = e^{-rT}\left[K\,N(-d_2) - F\,N(-d_1)\right],\qquad
  d_1 = \frac{\ln(F/K) + \sigma^2T/2}{\sigma\sqrt{T}},\quad d_2 = d_1 - \sigma\sqrt{T}.
\end{equation}

\begin{table}[t]\centering\small
\caption{E2: Parámetros y valorización del collar sobre WTI al \FechaValorizacion}\label{tab:e2-param}
\renewcommand{\arraystretch}{1.0}
\newcommand{\grupo}[1]{\rowcolor{gray!12}\multicolumn{3}{l}{\textit{#1}}\\}
\begin{tabularx}{\linewidth}{>{\raggedright\arraybackslash}p{5.3cm} r >{\raggedright\arraybackslash}X}
\toprule
\textbf{Parámetro} & \textbf{Valor} & \textbf{Fuente / supuesto} \\
\midrule
\grupo{Insumos de mercado}
Futuros $F$ nov / dic / ene (US\$/bl) & \WTIFutUno\ / \WTIFutDos\ / \WTIFut & NYMEX CL, 28-sep-2026 \\
Volatilidad implícita, $\sigma$ & \VolCrudo & Índice OVX; realizada a un año: \VolRealCL \\
Plazo al vencimiento, $T$ (días) & \PlazoDiasCapas & NYMEX: 7 días hábiles antes del 26 del mes previo \\
Tasa libre de riesgo, $r$ & \RUSDCont & SOFR, capitalización continua \\
\grupo{Estructura del collar (costo cero, por capa nov / dic / ene)}
Volumen cubierto: total / en collar (MMbl) & \InvCubiertoMMbl\ / \InvCollarMMbl & \CoberturaCrudoPct\ del inventario; \ContratosEquiv\ contratos (1:1) \\
Put comprado: strike (US\$/bl) & \CollarPutsCapas & \PutPctFwd\ de $F$ \\
Put comprado: prima media (US\$/bl) & \PrimaPut & \textcite{black1976}; delta medio $-$\DeltaPutAbs \\
Call vendido: strike (US\$/bl) & \CollarCallsCapas & Su prima iguala la del put: prima neta cero \\
\grupo{Resultado del inventario total si el WTI cae 30\,\% (US\$ MM)}
Sin cobertura & (\PerdidaSinCobTreinta) & Valor del inventario a mercado \\
Solo collar & (\PerdidaMaxCollar) & Piso del put sobre lo cubierto \\
\textbf{Propuesta: swap + collar} & \textbf{(\PerdidaMaxMixto)} & Solo put sobre el \CoberturaCrudoPct: US\$~\CostoPutTotal\ MM \\
\bottomrule
\end{tabularx}
\fuente{\textcite{petroperu2026eeff}, Nota 10; \textcite{yahoo2026cl}; \textcite{fred2026}; \textcite{bluegamma2026}. Elaboración propia con el modelo de \textcite{black1976}.}
\end{table}

% ARCHIVO GENERADO — NO EDITAR A MANO
\begin{table}[htbp]\centering\small
\caption{E2: Resultado sobre el inventario de crudo al vencimiento (US\$ millones)}\label{tab:e2-escenarios}
\begin{tabular}{lrrrrrrrr}
\toprule
 & \textbf{-30\%} & \textbf{-20\%} & \textbf{-10\%} & \textbf{$=F$} & \textbf{0\%} & \textbf{+10\%} & \textbf{+20\%} & \textbf{+30\%} \\
\midrule
WTI (US\$/bbl) & 64.8 & 74.1 & 83.3 & 86.5 & 92.6 & 101.9 & 111.1 & 120.4 \\
Sin cobertura & (80.6) & (53.7) & (26.9) & (17.6) & 0.0 & 26.9 & 53.7 & 80.6 \\
Swap (venta a 89.38) & (9.3) & (9.3) & (9.3) & (9.3) & (9.3) & (9.3) & (9.3) & (9.3) \\
Collar 77.90 / 98.08 & (42.6) & (42.6) & (26.9) & (17.6) & 0.0 & 15.9 & 15.9 & 15.9 \\
Solo put (prima pagada) & (57.0) & (57.0) & (41.2) & (31.9) & (14.4) & 12.5 & 39.4 & 66.2 \\
Propuesta: 50\% swap + 50\% collar & (26.0) & (26.0) & (18.1) & (13.5) & (4.7) & 3.3 & 3.3 & 3.3 \\
\bottomrule
\end{tabular}
\fuente{PETROPERÚ, EEFF auditados al 31.12.2025 (Gaveglio, Aparicio y Asociados, 2026), Nota 10; NYMEX CL vía Yahoo Finance; CBOE OVX vía FRED (28-sep-2026). Elaboración propia. Variación del WTI respecto del precio del 28-sep-2026. Signo negativo = pérdida.}
\end{table}


\begin{figure}[t]\centering
  \includegraphics[width=0.8\linewidth]{f_e2_payoff}
  \caption{E2: Perfil de resultados de la cobertura del inventario de crudo}\label{fig:e2}
  \fuente{\textcite{petroperu2026eeff}, Nota 10; \textcite{yahoo2026cl}; \textcite{fred2026}. Elaboración propia.}
\end{figure}

\textbf{Escenarios.} El mercado presenta una fuerte \emph{backwardation}: el futuro de enero está \Backwardation\
por debajo del contrato frente. El modelo de costo de acarreo, $F=S\,e^{(r-y)T}$, lo explica con un rendimiento de
conveniencia neto implícito de \ConvenienciaImplicita\ anual entre ambos contratos, muy superior a la tasa libre de
riesgo (\RUSDCont): en plena crisis de oferta el mercado paga por tener crudo disponible hoy \parencite{hull2022}.
Por eso toda venta a plazo fija un precio menor que el actual: sobre su volumen, el swap elimina la variabilidad a
cambio de US\$~\CostoSwap\ millones menos que el precio de hoy, pero su costo esperado frente al forward, que el mercado
ya descuenta, es nulo. Pero si el WTI converge al futuro (columna
$=F$ del Cuadro~\ref{tab:e2-escenarios}), el inventario sin cobertura pierde US\$~\PerdidaSinCobForward\ millones, más
que el swap. Si el WTI cae 30\,\%, el inventario pierde US\$~\PerdidaSinCobTreinta\ millones sin cobertura y
US\$~\PerdidaMaxCollar\ millones con el collar, que conserva la ganancia hasta el strike del call de cada capa y no
requiere una prima que la caja no puede financiar; comprar solo el put sobre todo lo cubierto
costaría US\$~\CostoPutTotal\ millones. En ese escenario la combinación propuesta reduce la pérdida a
US\$~\PerdidaMaxMixto\ millones. Si el WTI sube, la Compañía renuncia a parte de la ganancia (con un alza de
30\,\% gana US\$~\GananciaMaxMixto\ millones); aun así conviene, porque protege la caja de una empresa sin margen para absorber otra caída como la de 2025
\parencite{froot1993}. El OVX coincide con la volatilidad realizada del CL a un año (\VolRealCL; el \VolHistWTI\ de la sección~2 usa
promedios mensuales); con $\sigma=$~\VolSensBaja\ el call de enero apenas cambiaría (US\$~\CollarCallVolBaja), pero el
\emph{skew}, no modelado, obliga a confirmar los strikes con bancos.

El principal riesgo residual es el de base: PETROPERÚ compra crudos referidos al Brent, cuyo diferencial con el WTI
es hoy de US\$~\BrentWTI/bl, y crudos pesados con descuento. El ratio de mínima varianza $h^*=\rho\,\sigma_S/\sigma_F$
del Brent frente al futuro CL es \RatioHBrent, con una efectividad $\rho^2$ de \EfectividadBrent\ (\SemanasH\
semanas; para el WTI de Cushing: \RatioH\ y \EfectividadWTI) \parencite{hull2022}. Si la partida fuera el Brent, el
nocional subiría a \ContratosHBrent\ contratos y una quinta parte de la varianza quedaría sin cubrir. Se mantiene el
WTI porque los precios de venta locales siguen la paridad de importación sobre el WTI más spreads (Nota~11(viii),
p.~70), lo que lo hace un componente de riesgo identificable (NIIF~9, 6.3.7); los swaps sobre Brent (ICE) son la
alternativa si la serie interna de costos de compra se ajusta mejor a ese marcador.

\medskip\noindent
\begin{minipage}[t]{0.5\linewidth}\footnotesize\setlength{\tabcolsep}{2.5pt}% ARCHIVO GENERADO — NO EDITAR A MANO
\centering
\captionof{table}{E2: posiciones en barriles (MM)}\label{tab:e2-posiciones}
\begin{tabular}{lrr}
\toprule
 & \textbf{MMbl} & \textbf{WTI $-10$\,\% (US\$)} \\
\midrule
Físico: inventario de crudo & 2.90 & (26.86) \\
Swap (vende a precio fijo) & (1.16) & 10.75 \\
Collar (equivalente delta) & (0.62) & 5.75 \\
Global (físico + derivados) & 1.12 & (10.37) \\
\bottomrule
\end{tabular}
\fuente{\textcite{petroperu2026eeff}, Nota 10; NYMEX CL vía Yahoo Finance; CBOE OVX vía FRED (28-sep-2026). Elaboración propia. Positivo = largo en crudo. Collar: delta medio de las capas. Efecto lineal de una caída de 10 \% del WTI.}
\end{minipage}\hfill
\begin{minipage}[t]{0.47\linewidth}\textbf{Demostración de la cobertura.} Sobre su volumen (\InvSwapMMbl\ MMbl),
inventario y swap suman $-$US\$~\CostoSwap\ millones en cualquier escenario: lo que pierde el crudo lo gana el swap.
En barriles, los derivados reducen la posición larga de \InvCrudoMMbl\ a \PosGlobalMMbl\ MMbl equivalentes
(\PosGlobalPct): queda el 20\,\% abierto y la zona entre strikes del collar.\end{minipage}
\medskip

\textbf{Registro contable.} El swap y el collar se pactan a valor cero. Ante una caída instantánea de 10\,\% de la
curva de futuros, el collar sería un activo de US\$~\CollarValMenosDiez\ millones y el swap uno de
US\$~\SwapValDiez\ millones; con un alza de 10\,\%, serían pasivos de US\$~\CollarValMasDiez\ y \SwapValDiez\
millones. Cada capa se designa como cobertura de valor razonable del inventario de su mes: el ajuste por el riesgo
cubierto modifica el costo del inventario y va a resultados junto con el derivado. Al venderse la capa, la relación se
discontinúa (NIIF~9, 6.5.6) y se designa la siguiente. El Cuadro~\ref{tab:e2-escenarios} es económico: la pérdida
contable depende del costo y del valor neto de realización (NIC~2), no del spot \parencite{ifrs9}.
